\documentclass[12pt]{article}
\usepackage[margin=0.72in]{geometry}
\usepackage{lmodern}
\usepackage{microtype}
\usepackage{multicol}
\usepackage{fancyhdr}
\usepackage{titlesec}
\usepackage{enumitem}
\usepackage{xcolor}
\usepackage[hidelinks]{hyperref}

\setlength{\columnsep}{0.30in}
\setlength{\parindent}{1.05em}
\setlength{\parskip}{0.24em}
\setlength{\headheight}{15pt}
\titleformat{\section}{\large\bfseries\scshape}{}{0pt}{}
\titleformat{\subsection}{\normalsize\bfseries}{}{0pt}{}
\pagestyle{fancy}
\fancyhf{}
\fancyhead[L]{Augustine's Divided Will}
\fancyhead[R]{Reading Handout}
\fancyfoot[C]{\thepage}
\renewcommand{\headrulewidth}{0.4pt}

\newcommand{\focusbox}[1]{%
\vspace{0.4em}\noindent\colorbox{gray!12}{\begin{minipage}{0.96\linewidth}#1\end{minipage}}\vspace{0.5em}}

\begin{document}
\begin{center}
{\LARGE\bfseries Do We Love Shadows?}\par\vspace{0.25em}
{\large\scshape Reading 4: Augustine and the Divided Will}\par\vspace{0.45em}
\rule{0.82\textwidth}{0.4pt}\par\vspace{0.7em}
\end{center}

\section*{Vocabulary Preview}
\begin{multicols}{2}
\begin{itemize}[leftmargin=*, itemsep=0.18em]
\item \textbf{Confession}: An honest admission of truth before God and oneself.
\item \textbf{Conversion}: A turning of the soul toward God, truth, or a new way of life.
\item \textbf{Continence}: Self-control, especially over bodily desire.
\item \textbf{Concupiscence}: Strong disordered desire, especially desire that pulls against the good.
\item \textbf{Vanity}: Empty pride in appearances, reputation, or false success.
\item \textbf{Contrition}: Deep sorrow for sin or wrongdoing.
\item \textbf{Iniquity}: Serious moral wrong or sin.
\item \textbf{Will}: The power of choosing, desiring, and committing oneself to action.
\item \textbf{Habit}: A repeated action or desire that becomes difficult to break.
\item \textbf{Serenity}: Deep calm or inward peace.
\item \textbf{Admonition}: A warning, correction, or instruction.
\item \textbf{Alypius}: Augustine's close friend, present during the crisis in the garden.
\end{itemize}
\end{multicols}

\section*{Introduction}
Augustine lived from A.D. 354 to 430. In \textit{Confessions}, he tells the story of his restless search for truth and happiness. Plato gives the image of people chained among shadows. Homer shows that a dangerous illusion may be beautiful enough to desire. Bacon names the idols that distort the mind before it even begins to reason. Augustine now presses deeper: \textbf{what if a person sees the truth and still does not want to surrender what keeps him in darkness?}

In Book VIII, Augustine describes the crisis just before his conversion. He has already rejected many false ideas. He knows he must change. Yet he feels divided inside himself. One part of him wants healing; another part clings to old pleasures, habits, and self-protection. As you read, ask: \textbf{What if the strongest shadow is not ignorance, but a love or habit the soul refuses to give up?}

\focusbox{\textbf{What you need to know before this passage:} Augustine has spent years searching for wisdom, first through ambition and rhetoric, then through false religious systems, and finally through Christian teaching in Milan. A visitor named Pontitianus has just told Augustine about people who gave themselves wholly to God. Their example forces Augustine to face himself. He can no longer pretend that uncertainty is the reason for his delay. The real conflict is inside his own will.}

\focusbox{\textbf{Source note:} Adapted and modernized from Augustine, \textit{Confessions}, Book VIII, chapters 7, 9, 11, and 12, in E. B. Pusey's public-domain translation. This handout uses selected passages from a long work: a brief context summary is followed by the focused garden crisis, preserving Augustine's sequence, inward conflict, images, and theological language while using modern student-facing English.}

\begin{multicols}{2}
\section*{Reading}

\subsection*{1. Turned Toward Himself}
While Pontitianus was speaking, O Lord, you turned me back toward myself. I had placed myself behind my own back, unwilling to look at myself. But you set me before my own face, so that I could see how crooked, stained, and diseased I was.

I looked, and I was horrified. I did not know where to flee from myself. If I tried to turn my eyes away, Pontitianus kept telling his story, and you again set me against myself. You thrust me before my own eyes so that I might discover my sin and hate it. I had known it before, but I acted as if I did not know it. I winked at it and forgot it.

The more I loved the healthy souls I was hearing about, people who had given themselves wholly to you to be healed, the more I hated myself by comparison. Many years had passed since I first read Cicero's \textit{Hortensius} and was stirred to love wisdom. Yet I was still delaying. I still would not reject earthly happiness and devote myself to the search for truth.

In my youth I had prayed to you for self-control, but I had said, ``Grant me chastity and continence, but not yet.'' I was afraid you might hear me too soon and deliver me too soon from the disease of desire. I wanted that desire to be satisfied, not extinguished.

Now the day had come when I was laid bare to myself. My conscience began to rebuke me: ``Where are you now, tongue? You used to say that you would not cast off vanity because the truth was uncertain. But now the truth is certain, and still the burden holds you down.''

Inside, I was consumed with shame. I lashed my soul with rebukes, trying to make it follow me as I tried to follow you. But my soul drew back. It refused. All its arguments had been exhausted and defeated. Only a silent trembling remained. It feared, as if it were death, to be restrained from the stream of habit by which it was wasting away even to death.

\subsection*{2. The Mind Commands Itself}
Where does this strange thing come from? The mind commands the body, and the body obeys at once. The mind commands the hand to move, and the hand moves so quickly that command and obedience almost seem the same.

Yet when the mind commands itself to will, it is resisted. The mind commands the mind to will, and still the thing is not done. How can this be?

The answer is that the will is not whole. It partly wills and partly does not will. It commands only so far as it wills; and what it commands is not done so far as it does not will. If the will were entire, it would not need to command itself to will, because it would already will.

So it is not strange that a person can partly will and partly refuse. It is the sickness of the mind: it does not rise wholly, because it is lifted by truth but weighed down by habit. There are, in a sense, two wills, because one will is not whole. One part has what the other lacks.

\subsection*{3. Almost Free, But Still Held}
I was sick and tormented. I accused myself more severely than usual. I tossed and turned in my chain until it was almost broken. Only a slight link still held me, but it held me still.

You, O Lord, pressed upon me inwardly with severe mercy. You doubled the lashes of fear and shame, so that I would not give way again. If that last slender tie were not broken, it might recover strength and bind me more tightly than before.

I said within myself, ``Let it be done now. Let it be done now.'' As I spoke, I almost resolved. I almost did it, yet I did not do it. I did not fall all the way back to my old life, but I stopped nearby and took breath. I tried again. I came a little closer, then closer still. I almost touched and grasped it. Yet I still did not touch or grasp it.

I hesitated to die to death and to live to life. The worse path, to which I was accustomed, had more power over me than the better path, which I had not yet tried. The nearer the moment came when I would become another man, the more horror it struck into me. It did not drive me backward, but it held me suspended.

\subsection*{4. The Old Loves Whisper}
The toys of toys, the vanities of vanities, my old mistresses, still held me. They plucked at my garment of flesh and whispered softly, ``Are you casting us off? From this moment shall we never be with you again? From this moment will this or that never again be allowed to you?''

What were they suggesting by ``this or that''? What shameful things did they suggest, O my God? Let your mercy turn such things away from the soul of your servant.

Now I heard them much less than half. They no longer stood openly before me, contradicting me. Instead, they muttered behind my back. They secretly tugged at me as I was leaving, trying to make me look back. Still, they delayed me. I hesitated to burst free and leap toward the place where I was being called. A violent habit said to me, ``Do you think you can live without these things?''

But now that voice spoke very faintly. On the side toward which I had turned my face, and toward which I trembled to go, I saw the chaste dignity of Continence. She was serene, joyful without being loose or foolish, honestly inviting me to come and doubt nothing. She stretched out holy hands to receive me, hands full of many good examples: young men and women, people of every age, widows, and lifelong virgins. In all of them Continence was not barren, but fruitful with children of joy.

She seemed to smile at me and say, ``Can you not do what these young men and women have done? Did they do it by themselves, or by the Lord their God? Why do you stand in your own strength, and so fail to stand? Cast yourself upon Him. Do not fear. He will not withdraw and let you fall. Cast yourself upon Him without fear. He will receive you and heal you.''

I blushed deeply, because I still heard the muttering of those old toys, and I hung in suspense. Continence seemed to speak again: ``Shut your ears against those unclean desires on the earth, so that they may die. They tell you of delights, but not like the law of the Lord your God.''

This controversy in my heart was nothing but self against self. Alypius sat close beside me in silence, waiting to see the outcome of my unusual emotion.

\subsection*{5. Under the Fig Tree}
Then a deep reflection drew together all my misery from the secret bottom of my soul and heaped it up before the sight of my heart. A mighty storm rose in me, bringing a mighty shower of tears.

So that I could pour out those tears freely, I rose from Alypius. Solitude seemed better for the business of weeping. I went far enough away that even his presence would not burden me. He stayed where we had been sitting, greatly astonished.

I threw myself down, somehow, under a certain fig tree. I gave full vent to my tears, and the streams of my eyes poured out as an acceptable sacrifice to you. I spoke many things to you, not exactly in these words, but with this meaning: ``And you, O Lord, how long? How long, Lord, will you be angry forever? Remember not our former sins.'' I felt that I was held by them.

I sent up these sorrowful words: ``How long, how long? Tomorrow, and tomorrow? Why not now? Why should there not be, this very hour, an end to my uncleanness?''

\subsection*{6. Take Up and Read}
I was speaking and weeping in the bitter contrition of my heart when I heard a voice from a neighboring house. It sounded like the voice of a boy or girl, I do not know which, chanting again and again: ``Take up and read. Take up and read.''

Immediately my expression changed. I began to think carefully whether children usually sang such words in any kind of game, but I could not remember ever hearing anything like it. So I checked the flood of my tears and rose up. I interpreted the voice as nothing other than a command from God to open the book and read the first passage I should find.

I had heard of Antony, who once came in while the Gospel was being read and received the words as if they were spoken directly to him: ``Go, sell all that you have and give to the poor, and you shall have treasure in heaven; and come, follow me.'' By that word, he was immediately converted to you.

So I eagerly returned to the place where Alypius was sitting, for there I had laid the volume of the Apostle when I got up. I seized it, opened it, and silently read the passage on which my eyes first fell:

``Not in rioting and drunkenness, not in sexual immorality and shamelessness, not in strife and envying; but put on the Lord Jesus Christ, and make no provision for the flesh, to fulfill its desires.''

I would read no further, nor did I need to. Instantly, as the sentence ended, a light of serenity was poured into my heart, and all the darkness of doubt vanished away.

Then I put my finger in the book, or used some other mark, and closed it. With a calm face, I told Alypius what had happened. He asked to see what I had read. I showed him, and he looked further than I had read. The next words were, ``Receive the one who is weak in faith.'' He applied this to himself and joined me without any troubled delay.

Then we went in to my mother. We told her, and she rejoiced. We explained how it happened, and she leaped for joy and blessed you, who are able to do more than we ask or think. She saw that you had given her more for me than she had long begged with tears.

\section*{Reading Focus}
\begin{enumerate}[leftmargin=*, itemsep=0.2em]
\item Underline one sentence showing that Augustine already knows the truth but avoids facing himself.
\item Circle one sentence showing that habit acts like a chain.
\item Put a heart in the margin beside one sentence about love, desire, or old pleasures.
\item Put a star beside the moment when Augustine's will begins to turn.
\item Box one sentence that best answers this question: What if the strongest shadow is a love or habit the soul refuses to give up?
\end{enumerate}

\section*{Loves Inventory}
Make two columns in your notebook.

\begin{itemize}[leftmargin=*, itemsep=0.2em]
\item \textbf{Lower loves:} List three things Augustine is tempted to keep loving in the wrong way.
\item \textbf{Better-ordered loves:} For each lower love, write what a healed or rightly ordered love would look like.
\end{itemize}

\section*{Inner-Conflict Dialogue}
Write a short dialogue between two voices inside Augustine.

\begin{itemize}[leftmargin=*, itemsep=0.2em]
\item \textbf{Voice 1} wants comfort, old pleasure, delay, or self-protection.
\item \textbf{Voice 2} wants truth, healing, freedom, and conversion.
\end{itemize}

Your dialogue should be at least ten lines long. Each voice should make a serious argument.

\section*{Claim Question}
Answer in one clear paragraph: \textbf{Why might a person resist truth even after recognizing it?} State a clear claim, use one specific moment from Augustine, explain how love or habit can keep someone attached to shadows, and connect Augustine to Plato, Homer, or Bacon.

\section*{Handout Summary}
Finish by writing a short summary of this handout. In three to five sentences, explain Augustine's inner conflict, what finally changes in the garden, and how this reading adds to Plato, Homer, and Bacon.

\end{multicols}
\end{document}
